% =====================================================================
\chapter{Специальные функции и голос}
\label{ch:special}
% =====================================================================
\chapterintro
  {какие действия может выполнять пульт по условию: принудительно задавать канал,
   говорить, вибрировать, писать логи, сбрасывать таймеры, менять громкость; как
   записать свои голосовые сообщения}
  {модель; для голоса --- голосовой пакет на SD-карте}

\section{Страница Special Functions}

\mpath{MDL\sep Special Functions} --- список правил вида «\emph{если} переключатель включён,
\emph{то} выполнить действие». Строки \sw{SF1}…\sw{SF64}. Такие же правила, но общие для всех
моделей, задаются в \mpath{SYS\sep Global Functions}.

\begin{center}
\lcd{\inv{SPECIAL FUNCTIONS}\\
SF1\quad SD↓\quad Override CH3\quad -100\quad ☑\\
SF2\quad SA↓\quad Play track\quad arm\\
SF3\quad L10\quad Play track\quad batlow\quad 10s\\
SF4\quad ON\quad\ \ \ SD logs\quad 0.2s}
\end{center}

У каждой строки:
\begin{itemize}
  \item \menu{Switch} --- условие: физический или логический переключатель, \sw{ON}
        (всегда), \sw{One} (один раз при загрузке модели), состояния телеметрии и др.;
  \item \menu{Function} --- действие;
  \item параметры действия;
  \item \menu{Repeat} --- для звуков: проиграть один раз (\code{1x}), при каждом включении
        или повторять каждые N секунд;
  \item \menu{Enable} --- включить/выключить строку (флажок).
\end{itemize}

\section{Список функций}

\begin{center}\small
\renewcommand{\arraystretch}{1.15}
\begin{tabularx}{\linewidth}{@{}p{34mm}X@{}}
\toprule
Функция & Что делает \\
\midrule
\menu{Override CHx} & принудительно задаёт значение канала (отсечка газа, «замок» арма) \\
\menu{Trainer} & передаёт управление ученику (все оси или отдельные) \\
\menu{Instant trim} & переносит текущее положение стиков в триммеры \\
\menu{Reset} & сбрасывает таймер, телеметрию (мин./макс.) или всё сразу \\
\menu{Set timer} & устанавливает значение таймера \\
\menu{Adjust GVx} & меняет глобальную переменную: на число, на значение источника,
  $+1$/$-1$ \\
\menu{Volume} & регулирует громкость источником (например, крутилкой) \\
\menu{Set failsafe} & записывает текущие положения каналов как failsafe (для модулей,
  где это поддерживается) \\
\menu{Range check}, \menu{Bind} & запуск проверки дальности и привязки модуля \\
\menu{Play sound} & системный звук (писк, трель и т.\,п.) \\
\menu{Play track} & воспроизводит WAV-файл из папки голосового пакета \\
\menu{Play value} & голосом называет значение источника (напряжение, таймер, RSSI) \\
\menu{Lua script} & запускает функциональный Lua-скрипт (глава~\ref{ch:lua}) \\
\menu{Background music} & фоновая музыка (WAV) \\
\menu{Vario} & звуковой вариометр --- тон зависит от вертикальной скорости \\
\menu{Haptic} & вибрация \\
\menu{SD logs} & запись телеметрии и каналов в CSV-файл на SD-карту \\
\menu{Backlight} & яркость подсветки от источника \\
\menu{Screenshot} & снимок экрана в папку \code{SCREENSHOTS} \\
\bottomrule
\end{tabularx}
\end{center}

\section{Override: принудительное значение канала}

\begin{recipe}{отсечка газа}
\menu{Switch} = \sw{SD↓}, \menu{Function} = \menu{Override CH3}, значение $-100$,
флажок \menu{Enable} включён. Пока \sw{SD} внизу, канал газа равен $-100$ независимо
от стика. Добавьте \sw{SD} в предполётные проверки (положение «отсечка включена»),
и модель не запустит мотор при включении.
\end{recipe}

\begin{caution}
Для квадрокоптеров с Betaflight основная защита --- \emph{арм} на \sw{CH5}, а не отсечка
газа. Отсечку через \menu{Override} для квадрокоптеров обычно не делают: полётный контроллер
в арме держит минимальные обороты, и принудительный газ $-100$ не отключит моторы.
\end{caution}

\section{Голосовые сообщения}

\subsection{Готовые фразы и значения}

\begin{itemize}
  \item \menu{Play value} с источником \sw{Tmr1} --- «три минуты двадцать секунд»;
  \item \menu{Play value} с источником \sw{RxBt} --- «пятнадцать и две десятых вольта»;
  \item \menu{Play value} с источником \sw{RQly} --- «девяносто восемь процентов».
\end{itemize}

\begin{recipe}{голос по кнопке}
\menu{Switch} = \sw{SF↓} (кнопка без фиксации), \menu{Play value} \sw{RxBt}, \menu{Repeat} =
\code{1x}. Нажали кнопку --- пульт сказал напряжение батареи модели.
\end{recipe}

\subsection{Озвучка переключателей}

Для каждого положения переключателя можно проигрывать файл: \sw{SA↓} → \code{armed},
\sw{SA↑} → \code{disarm}. В стандартном голосовом пакете уже есть много фраз
(«armed», «disarmed», «flight mode», «angle», «acro», «beeper» и т.\,п.). Список файлов --- в
папке \code{SOUNDS/<язык>} на SD-карте.

\subsection{Свои звуковые файлы}

\begin{enumerate}
  \item Запишите или синтезируйте фразу и сохраните её как \textbf{WAV, 16 бит, моно, PCM}
        с частотой дискретизации 32\,кГц (подходят также 16 и 8\,кГц).
  \item Дайте файлу короткое латинское имя без пробелов (не длиннее 6--8 символов),
        например \code{batlow.wav}.
  \item Положите в \code{SOUNDS/ru} (или папку вашего языка).
  \item В специальной функции выберите \menu{Play track} → \code{batlow}.
\end{enumerate}

Для синтеза удобно использовать любой сервис «текст в речь» и аудиоредактор (например,
Audacity) для конвертации в нужный формат.

\section{Запись логов}

\begin{recipe}{логи полёта}
\menu{Switch} = \sw{SA↓} (арм) или \sw{ON}, \menu{Function} = \menu{SD logs}, период 0,2\,с.
Во время полёта в папку \code{LOGS} пишется файл \code{ИмяМодели-ГГГГ-ММ-ДД.csv}: время,
все датчики телеметрии, положения стиков и переключателей. Файл открывается в табличном
редакторе или в EdgeTX Companion; координаты GPS можно наложить на карту.
\end{recipe}

\section{Global Functions}

Функции, полезные для \emph{любой} модели, удобнее задать один раз в
\mpath{SYS\sep Global Functions}:
\begin{itemize}
  \item \menu{Volume} на крутилке (если она свободна во всех моделях);
  \item \menu{Screenshot} по долгому нажатию кнопки (для инструкций и отчётов);
  \item \menu{Play value} напряжения батареи пульта по кнопке.
\end{itemize}
Модель использует глобальные функции, только если в её \menu{Setup} включено
\menu{Use global functions}.

\begin{summary}
\begin{itemize}
  \item Специальная функция = условие + действие (+ повтор, + флажок).
  \item \menu{Override} --- безусловное значение канала; \menu{Play track}/\menu{Play value} ---
        голос; \menu{SD logs} --- логи.
  \item Свои звуки --- WAV 16 бит моно в \code{SOUNDS/<язык>}.
\end{itemize}
\end{summary}
